%! Author = chtompki
%! Date = 1/14/26
\documentclass[12pt]{amsart}
% Default margins are too wide all the way around. I reset them here
\setlength{\hoffset}{-.75in}
\setlength{\textwidth}{6.5in}
\setlength{\voffset}{-.5in}
\setlength{\textheight}{9.0in}
\usepackage{amsmath}
\usepackage{amssymb}
\usepackage{amsthm}
\usepackage{enumitem}
\usepackage{setspace}
\usepackage{graphicx}
\usepackage[colorlinks=true, urlcolor=black, linkcolor=black]{hyperref}

\newenvironment{solution}
{\begin{proof}[Solution]}
{\end{proof}}

\title{Numerical Analysis\\Guided Lab 3}
\author{Rob Tompkins}
\date{\today}

\begin{document}
    \maketitle
    \date{\today}
    \begin{onehalfspace}
        Note, all the files associated with Guided Lab 3 are checked into the : \\
        \underline{https://github.com/chtompki/MathSpring2026/tree/main/NumericalAnalysis/GuidedLab3}

        Problems 1, 2, and 3, were predominantly only writing code. The code can be found at:
        \underline{\href{https://github.com/chtompki/MathSpring2026/tree/main/NumericalAnalysis/GuidedLab3/naturalCubicSplines.m}{naturalCubicSplines.m}}.
        Though, the code for these three problems also follows as
        \begin{tiny}
        \begin{verbatim}
function [A,rvec,c,a,b,d] = naturalCubicSplines(x,y)
    if length(x) < 2 | length(y) < 2
        error('Input vectors are not long enough, please make them of length two or more and be of the same size')
    end
    if length(x) ~= length(y)
        error('Input vectors are not the same length')
    end
    A=zeros(length(x));
    rvec=zeros(length(x),1);
    h=zeros(1,length(x));

    % Begin Problem 1
    A(1,1) = 1;
    A(length(x),length(x)) = 1;

    h(1) = x(2) - x(1);

    for i = 2:length(x)-1
        h(i) = x(i+1) - x(i);
        A(i,i-1) = h(i-1);
        A(i,i) = 2*(h(i-1)+h(i));
        A(i,i+1) = h(i);
        rvec(i) = 3*((y(i+1)-y(i))/h(i) - (y(i)-y(i-1))/h(i-1));
    end
    % End Problem 1

    % Problem 2 - Solve the system of equations
    c = A \ rvec;

    % Problem 3 - Calculate the coefficients a, b, and d
    a = zeros(length(x),1);
    for i=1:length(x)
        a(i) = y(i);
    end

    b = zeros(length(x)-1, 1);
    d = zeros(length(x)-1, 1);

    for i = 1:length(x)-1
        b(i) = (c(i+1) - c(i))/(3*h(i));
        d(i) = (a(i+1) - a(i))/h(i) - (h(i)*(2*c(i) + c(i+1)))/3;
    end
end
        \end{verbatim}
        \end{tiny}
    \end{onehalfspace}
\end{document}
